\clearpage
\fancyhead[L]{\small INSTRUMENTE.}
\SourceLeaf{440}
\begin{center}
{\large INSTRUMENTE.}
\end{center}

\begin{center}
\noindent\rule{\linewidth}{0.4pt}

Göttingische gelehrte Anzeigen. Stück 187. Seite 1865..1866. 1820. November 20.

\noindent\rule{\linewidth}{0.4pt}
\end{center}

Se Königl. Hoheit der Herzog von \textsc{Clarence} hat der Universität einen Beweis
seines gnädigen Andenkens gegeben, indem Er die Bibliothek mit einer kostbaren Sammlung
von Seecharten in 182 Blättern beschenkt hat. Es sind dies die in London in dem
Hydrographical Office erscheinenden, und mit dessen Stempel bezeichneten Charten; die
nicht in den Buchhandel kommen, sondern nur für den Gebrauch der Königl. Marine bestimmt
sind. Die Sammlung umfasst nicht bloss die Europäischen Gewässer, sondern auch den
grössten Theil der Küsten von Africa, America, Ost- und Westindien, von denen der
Herzog, selber Seemann, viele aus eigner Ansicht kennt. Mehr wird es nicht bedürfen,
sowohl ihren Werth als unsre Dankbarkeit zu bezeichnen, da ein Verzeichniss der
Einzelnen dem Zweck dieser Blätter nicht angemessen sein würde.

Auch von Sr. Königl. Hoheit dem Herzog von \textsc{Sussex} erhielt die Universität einen
Beweis seiner huldreichen Erinnerung, durch zwei Geschenke, welche Höchstdieselben der
Sternwarte gemacht haben. Es bestehen diese in einer von \textsc{Hardy} in London
verfertigten Tertienuhr, von einer besondern Einrichtung und vorzüglich schöner Arbeit,
und in einem Apparat, den man ein verkehrtes Pendel nennen könnte, und wodurch die
Anwendbarkeit einer sinnreichen Idee ins Licht gesetzt wird. So wie das gewöhnliche
Pendel oben befestigt ist und an der Feder \emph{hängt}, so ist dieses unten fest, und
\emph{steht} auf der verhältnissmässig etwas starken Feder. Die Pendelstange ist ein 5
Zoll langer Cylinder, auf welchem oben ein kleines kugelförmiges Gewicht aufgeschraubt
ist. Oben endigt sich die Stange in eine feine Spitze, die auf einem in 20 Theile
getheilten, Einen Zoll grossen Gradbogen spielt. Jeder dieser Theile beträgt also
eigentlich etwas über 34 Minuten, allein die auf der Sternwarte gemachten Versuche haben
ergeben, dass schon eine Neigung des Apparats von 57 Secunden einen Ausschlag von einem
Theile gibt, oder dass dieses Pendel eine eben so grosse Empfindlichkeit hat, wie ein
gewöhnliches von 15 Fuss Länge. Da das Kügelchen höher und tiefer geschraubt werden
kann, so kann man, wenn man will, es leicht so abgleichen, dass der Ausschlag Eines
Theils genau Einer Minute entspreche. Man sieht hieraus, dass dieser Apparat, mit den
nöthigen Correctionsschrauben versehen, und mit gehöriger Vorsicht angewandt, in manchen
Fällen mit Vortheil die Stelle einer Libelle oder eines Loths vertreten kann.

\clearpage
\fancyhead[L]{\small COMET.}
\hypertarget{native-440-part-lower}{}
\begin{center}
{\large COMET.}
\end{center}

\begin{center}
\noindent\rule{\linewidth}{0.4pt}

Göttingische gelehrte Anzeigen. Stück 23. Seite 217. 1821 Februar 10.

\noindent\rule{\linewidth}{0.4pt}
\end{center}

Der von Herrn \textsc{Nicolet} in Paris am 21. Januar d. J. entdeckte Comet wurde auf
der hiesigen Sternwarte am 30. Januar von Hrn. Hofr. \textsc{Gauss} aufgefunden und
beobachtet. Um $7^{\mathrm{u}}\,34^{\mathrm{m}}\,32^{\mathrm{s}}$ Mittlere Zeit war des
Cometen

\begin{center}
\begin{tabular}{ll}
Gerade Aufsteigung . . . . . . . . & $359^\circ\;27'\;7''$ \\
Die nördliche Abweichung . . . . . . & $16\quad4\quad36$ \\
\end{tabular}
\end{center}

Dieser Comet ist dem blossen Auge unsichtbar, aber schon im Cometensucher gut zu
erkennen, wo sich auch der Schweif in der Länge von einigen Graden zeigt. Von einem
begrenzten Kern ist keine Spur.
